\documentclass[10pt,a4paper]{amsart}
\usepackage{amsmath,amssymb,amsthm,graphicx}
% For commutative diagrams we use Xy-pic
\usepackage[all]{xy}

% We use 2cell for 2-commutative diagrams.
\xyoption{2cell}
% Theorem environments.
%
\theoremstyle{plain}
\newtheorem{theorem}[subsection]{Theorem}
\newtheorem{proposition}[subsection]{Proposition}
\newtheorem{lemma}[subsection]{Lemma}

\theoremstyle{definition}
\newtheorem{definition}[subsection]{Definition}
\newtheorem{example}[subsection]{Example}
\newtheorem{exercise}[subsection]{Exercise}
\newtheorem{situation}[subsection]{Situation}

\theoremstyle{remark}
\newtheorem{remark}[subsection]{Remark}
\newtheorem{remarks}[subsection]{Remarks}

\numberwithin{equation}{subsection}

% Macros
%
\def\lim{\mathop{\mathrm{lim}}\nolimits}
\def\colim{\mathop{\mathrm{colim}}\nolimits}
\def\Spec{\mathop{\mathrm{Spec}}}
\def\Hom{\mathop{\mathrm{Hom}}\nolimits}
\def\Ext{\mathop{\mathrm{Ext}}\nolimits}
\def\SheafHom{\mathop{\mathcal{H}\!\mathit{om}}\nolimits}
\def\SheafExt{\mathop{\mathcal{E}\!\mathit{xt}}\nolimits}
\def\Sch{\mathit{Sch}}
\def\Mor{\mathop{\mathrm{Mor}}\nolimits}
\def\Ob{\mathop{\mathrm{Ob}}\nolimits}
\def\Sh{\mathop{\mathit{Sh}}\nolimits}
\def\NL{\mathop{N\!L}\nolimits}
\def\CH{\mathop{\mathrm{CH}}\nolimits}
\def\proetale{{pro\text{-}\acute{e}tale}}
\def\etale{{\acute{e}tale}}
\def\QCoh{\mathit{QCoh}}
\def\Ker{\mathop{\mathrm{Ker}}}
\def\Im{\mathop{\mathrm{Im}}}
\def\Coker{\mathop{\mathrm{Coker}}}
\def\Coim{\mathop{\mathrm{Coim}}}

% Boxtimes
%
\DeclareMathSymbol{\boxtimes}{\mathbin}{AMSa}{"02}

%
% Macros for moduli stacks/spaces
%
\def\QCohstack{\mathcal{QC}\!\mathit{oh}}
\def\Cohstack{\mathcal{C}\!\mathit{oh}}
\def\Spacesstack{\mathcal{S}\!\mathit{paces}}
\def\Quotfunctor{\mathrm{Quot}}
\def\Hilbfunctor{\mathrm{Hilb}}
\def\Curvesstack{\mathcal{C}\!\mathit{urves}}
\def\Polarizedstack{\mathcal{P}\!\mathit{olarized}}
\def\Complexesstack{\mathcal{C}\!\mathit{omplexes}}
% \Pic is the operator that assigns to X its picard group, usage \Pic(X)
% \Picardstack_{X/B} denotes the Picard stack of X over B
% \Picardfunctor_{X/B} denotes the Picard functor of X over B
\def\Pic{\mathop{\mathrm{Pic}}\nolimits}
\def\Picardstack{\mathcal{P}\!\mathit{ic}}
\def\Picardfunctor{\mathrm{Pic}}
\def\Deformationcategory{\mathcal{D}\!\mathit{ef}}
\newenvironment{reference}{}{}
\newenvironment{slogan}{}{}
\newenvironment{history}{}{}
\usepackage{cleveref}
\crefname{originallemma}{original lemma}{original lemmas}
\Crefname{originallemma}{Original lemma}{Original lemmas}
\makeatletter
\newlabel{lemma-recurring-symmetric-real}{{0C5X}{source}{}{sourcepointer-1}{}}
\newlabel{lemma-recurring-symmetric-real@cref}{{[originallemma][0][]0C5X}{[1][0][]source}{}{}{}}
\makeatother
\crefname{originallemma}{original lemma}{original lemmas}
\Crefname{originallemma}{Original lemma}{Original lemmas}
\makeatletter
\newlabel{lemma-orthogonal-direct-sum}{{0C6V}{source}{}{sourcepointer-2}{}}
\newlabel{lemma-orthogonal-direct-sum@cref}{{[originallemma][0][]0C6V}{[1][0][]source}{}{}{}}
\makeatother
\crefname{originallemma}{original lemma}{original lemmas}
\Crefname{originallemma}{Original lemma}{Original lemmas}
\makeatletter
\newlabel{lemma-coker}{{0C6W}{source}{}{sourcepointer-3}{}}
\newlabel{lemma-coker@cref}{{[originallemma][0][]0C6W}{[1][0][]source}{}{}{}}
\makeatother
\title{The order of the l-torsion of the cokernel of a recurring symmetric integer matrix is bounded by the number of its positive off-diagonal pairs}
\author{The Stacks Project Authors}
\date{stacks-project:models.tex:local-snapshot-20261012; GNU FDL 1.2 or later}
\begin{document}
\maketitle
\noindent\textit{Editorial note.} The counted unit is the Stacks Project Lemma lemma-recurring-symmetric-integer at source lines 517-533 together with its proof at source lines 535-628, a straight adjacent pairing in the chapter file. The chapter is the upstream file models.tex; the pool directory carries the longer name stacks-models and is not used as the chapter name anywhere in this card.\par\medskip
\noindent\textit{Editorial note.} This card has no context span: everything the statement and the proof use is either inside the counted span or supplied through the three navigation pointers below. In particular the matrix appears explicitly in the statement, so its normalisation inside the proof is self-contained.\par\medskip
\noindent\textit{Editorial note.} The retained proof names three results of the same chapter, the rank statement for recurring symmetric matrices, the duality statement for a lattice modulo an orthogonal direct sum, and the identification of the torsion of a cokernel with a dual quotient. All three are declared as non-load-bearing navigation pointers with their stable Stacks tags looked up in the shared tag index. The third one in particular has an omission inside its own proof, which is why this lane does not count it; pointing at its statement from here is nevertheless valid, and its interval is not reproduced.\par\medskip
\noindent\textit{Editorial note.} No citation occurs in the retained spans, so no bibliography entry is supplied.\par\medskip
\noindent\textit{Allowed standard prerequisites (editorial).}\begin{itemize}
\item Recurring symmetric integer matrices and the lattices attached to a graph of nonnegative entries. The lemma concerns a symmetric integer matrix whose off-diagonal entries are nonnegative, together with a positive integer vector in its kernel such that no proper nonempty subset of indices fails to interact with its complement; for such data one forms the free modules on the indices and on the pairs with positive entry, equips them with the diagonal pairings given by one and by the entry itself, and considers the incidence map and its adjoint. The chapter proves separately, for the matrix in question, that the quadratic form is negative semidefinite with kernel spanned by the weight vector, and it proves a duality statement for a lattice modulo an orthogonal direct sum. The author's own lemmas of this section are the allowed inputs.. The counted lemma is the arithmetic estimate that the chapter uses when it computes the torsion of the cokernel of the intersection matrix of a numerical type.
\end{itemize}\medskip
\ifdefined\hypertarget\hypertarget{sourcepointer-1}{}\fi
\noindent\textit{Source reference.} Original lemma 0C5X: focused-seeds/AGP/sources/stacks-snapshot-20261011/tags.txt, source lines 15277--15277. Navigation pointer invoked by the retained proof for the rank of the matrix under the recurrence conditions; that lemma is counted by a separate card of this lane, so its proof is not reproduced here. The number is the stable Stacks tag of the target, looked up in the shared tag index.\par\medskip
\ifdefined\hypertarget\hypertarget{sourcepointer-2}{}\fi
\noindent\textit{Source reference.} Original lemma 0C6V: focused-seeds/AGP/sources/stacks-snapshot-20261011/tags.txt, source lines 15310--15310. Navigation pointer invoked by the retained proof for the duality statement applied to the image of the incidence map and the kernel of its adjoint; that lemma is counted by a separate card of this lane, so its proof is not reproduced here. The number is the stable Stacks tag of the target, looked up in the shared tag index.\par\medskip
\ifdefined\hypertarget\hypertarget{sourcepointer-3}{}\fi
\noindent\textit{Source reference.} Original lemma 0C6W: focused-seeds/AGP/sources/stacks-snapshot-20261011/tags.txt, source lines 15311--15311. Navigation pointer invoked by the retained proof for the final identification of the torsion module of the cokernel with the dual quotient of the image; the cited statement is an allowed input of the same section and its proof is not reproduced here. The number is the stable Stacks tag of the target, looked up in the shared tag index.\par\medskip
\begin{lemma}
\label{lemma-recurring-symmetric-integer}
Let $A = (a_{ij})$ be a symmetric $n \times n$ integer matrix with
$a_{ij} \geq 0$ for $i \not = j$. Let $m = (m_1, \ldots, m_n)$ be an
integer vector with $m_i > 0$. Assume
\begin{enumerate}
\item $Am = 0$,
\item there is no proper nonempty subset $I \subset \{1, \ldots, n\}$
such that $a_{ij} = 0$ for $i \in I$ and $j \not \in I$.
\end{enumerate}
Let $e$ be the number of pairs $(i, j)$ with $i < j$ and $a_{ij} > 0$.
Then for $\ell$ a prime number coprime with all $a_{ij}$ and $m_i$
we have
$$
\dim_{\mathbf{F}_\ell}(\Coker(A)[\ell]) \leq 1 - n + e
$$
\end{lemma}

\begin{proof}
By Lemma \ref{lemma-recurring-symmetric-real} the rank of $A$ is $n - 1$.
The composition
$$
\mathbf{Z}^{\oplus n} \xrightarrow{\text{diag}(m_1, \ldots, m_n)}
\mathbf{Z}^{\oplus n} \xrightarrow{(a_{ij})}
\mathbf{Z}^{\oplus n} \xrightarrow{\text{diag}(m_1, \ldots, m_n)}
\mathbf{Z}^{\oplus n}
$$
has matrix $a_{ij}m_im_j$. Since the cokernel of the first and last
maps are torsion of order prime to $\ell$ by our restriction on $\ell$
we see that it suffices to prove the lemma for the matrix
with entries $a_{ij}m_im_j$. Thus we may assume $m = (1, \ldots, 1)$.

\medskip\noindent
Assume $m = (1, \ldots, 1)$. Set $V = \{1, \ldots, n\}$ and
$E = \{(i, j) \mid i < j\text{ and }a_{ij} > 0\}$. For
$e = (i, j) \in E$ set $a_e = a_{ij}$. Define maps
$s, t : E \to V$ by setting $s(i, j) = i$ and $t(i, j) = j$.
Set
$\mathbf{Z}(V) = \bigoplus_{i \in V} \mathbf{Z}i$ and
$\mathbf{Z}(E) = \bigoplus_{e \in E} \mathbf{Z}e$.
We define symmetric positive definite integer valued pairings
on $\mathbf{Z}(V)$ and $\mathbf{Z}(E)$ by setting
$$
\langle i, i \rangle = 1\text{ for }i \in V, \quad
\langle e, e \rangle = a_e\text{ for }e \in E
$$
and all other pairings zero. Consider the maps
$$
\text{d} : \mathbf{Z}(V) \to \mathbf{Z}(E), \quad
i \longmapsto
\sum\nolimits_{e \in E,\ s(e) = i} e - \sum\nolimits_{e \in E,\ t(e) = i} e
$$
and
$$
\text{d}^*(e) = a_e(s(e) - t(e))
$$
A computation shows that
$$
\langle d(x), y\rangle = \langle x, \text{d}^*(y) \rangle
$$
in other words, $\text{d}$ and $\text{d}^*$ are adjoint. Next we compute
\begin{align*}
\text{d}^*\text{d}(i)
& = 
\text{d}^*(
\sum\nolimits_{e \in E,\ s(e) = i} e - \sum\nolimits_{e \in E,\ t(e) = i} e) \\
& =
\sum\nolimits_{e \in E,\ s(e) = i} a_e(s(e) - t(e)) -
\sum\nolimits_{e \in E,\ t(e) = i} a_e(s(e) - t(e))
\end{align*}
The coefficient of $i$ in $\text{d}^*\text{d}(i)$ is
$$
\sum\nolimits_{e \in E,\ s(e) = i} a_e +
\sum\nolimits_{e \in E,\ t(e) = i} a_e = - a_{ii}
$$
because $\sum_j a_{ij} = 0$ and the coefficient of
$j \not = i$ in $\text{d}^*\text{d}(i)$ is $-a_{ij}$.
Hence $\Coker(A) = \Coker(\text{d}^*\text{d})$.

\medskip\noindent
Consider the inclusion
$$
\Im(\text{d}) \oplus \Ker(\text{d}^*) \subset \mathbf{Z}(E)
$$
The left hand side is an orthogonal direct sum. Clearly
$\mathbf{Z}(E)/\Ker(\text{d}^*)$ is torsion free.
We claim $\mathbf{Z}(E)/\Im(\text{d})$ is torsion free as well.
Namely, say $x = \sum x_e e \in \mathbf{Z}(E)$ and $a > 1$ are such
that $ax = \text{d}y$ for some $y = \sum y_i i \in \mathbf{Z}(V)$.
Then $a x_e = y_{s(e)} - y_{t(e)}$. By property (2) we conclude
that all $y_i$ have the same congruence class modulo $a$.
Hence we can write $y = a y' + (y_1, y_1, \ldots, y_1)$.
Since $\text{d}(y_1, y_1, \ldots, y_1) = 0$ we conclude
that $x = \text{d}(y')$ which is what we had to show.

\medskip\noindent
Hence we may apply Lemma \ref{lemma-orthogonal-direct-sum}
to get injective maps
$$
\Im(\text{d})^\#/\Im(\text{d}) \leftarrow
\mathbf{Z}(E)/(\Im(\text{d}) \oplus \Ker(\text{d}^*)) \rightarrow
\Ker(\text{d}^*)^\#/\Ker(\text{d}^*)
$$
whose cokernels are annihilated by the product of the $a_e$
(which is prime to $\ell$). Since $\Ker(\text{d}^*)$ is
a lattice of rank $1 - n + e$ we see that the proof is complete
if we prove that there exists an isomorphism
$$
\Phi : M_{torsion} \longrightarrow \Im(\text{d})^\#/\Im(\text{d})
$$
This is proved in Lemma \ref{lemma-coker}.
\end{proof}

\end{document}
